\documentclass[11pt]{article}
\usepackage[margin=1.1in]{geometry}
\usepackage{amsmath,amssymb,amsthm}
\usepackage[colorlinks=true,linkcolor=blue,citecolor=blue]{hyperref}
\usepackage{booktabs}
\usepackage{longtable}
\usepackage{enumitem}

\newcommand{\Oo}{\mathcal O}
\newcommand{\ch}{\operatorname{ch}}
\newcommand{\tr}{\operatorname{tr}}
\newcommand{\C}{\mathbb C}
\newcommand{\Z}{\mathbb Z}
\newcommand{\Q}{\mathbb Q}
\newcommand{\R}{\mathbb R}
\newcommand{\Cl}{\operatorname{cl}}
\newcommand{\CH}{\operatorname{CH}}
\newcommand{\End}{\operatorname{End}}
\newcommand{\Hdg}{\operatorname{Hdg}}
\newcommand{\Emb}{\operatorname{Emb}}
\newcommand{\Gal}{\operatorname{Gal}}
\newcommand{\trd}{\operatorname{trd}}
\newcommand{\Nrd}{\operatorname{Nrd}}
\newcommand{\Gr}{\operatorname{Gr}}
\newcommand{\Sym}{\operatorname{Sym}}
\newcommand{\MT}{\mathrm{MT}}
\newcommand{\NS}{\operatorname{NS}}
\newcommand{\Hom}{\operatorname{Hom}}
\newtheorem{theorem}{Theorem}[section]
\newtheorem{lemma}[theorem]{Lemma}
\newtheorem{corollary}[theorem]{Corollary}
\newtheorem{proposition}[theorem]{Proposition}
\newtheorem{remark}[theorem]{Remark}
\newtheorem{definition}[theorem]{Definition}

\title{The Albert Fold-Birth\\
\large Round 49: the true input space at every stratification,
the false hypothesis killed, and the first derived landing cycles}
\author{The campaign (rounds 1--48) and this round's strike}
\date{2026-11-08 (Round 49)}

\begin{document}
\maketitle

\begin{abstract}
Round 48 reduced the remaining falsifiable line for Q2 --- the Clay
statement, \emph{every rational $(p,p)$ class on every smooth
projective complex variety is algebraic} --- to the named pair
$(C1,C2)$: the landing chain (the supplied field
\texttt{FoldProcessOn.landing\_cycle}) and the instantiation
hypothesis (every variety's extracted input is a two-letter world).
This round kills the second, replaces it with the true stratification,
generalizes the fold equations to the whole arithmetic input space, and
breaks the supplied field's monopoly at two seeds.
\textbf{Theorem 1 (the death certificate)}: $(C2)$ is refuted as a
blanket modeling claim --- the extraction is \emph{empty} at the
flagship open input (the cubic fourfold: $H^3=0$), and a non-empty
letterless input exists (the very general product of non-isogenous
abelian surfaces: $M_{\mathrm{mid}}\neq 0$ with
$\End_{\Hdg}(M_{\mathrm{mid}})=\Q\oplus\Q$, which admits no quadratic
letter). \textbf{Theorem 2 (the Albert stratification)}: the true world
model --- every input with $M_{\mathrm{mid}}\neq 0$ lands in exactly
one Albert cell of $\End^0(J_{\mathrm{mid}})$ (the trivial $\Q$-residue
or a simple component of Albert type I--IV), and the R47/R48 quadratic
strata are exactly the degree-$2$ cells.
\textbf{Theorem 3 (the orbit law)}: the phase-pair cancellation --- the
heart of the five fold equations --- generalizes from the quadratic
case to \emph{every} number field of even degree $n=2f$: each
conjugate pair of eigencells contributes the eigenvalue identity
$(\varepsilon i)(-\varepsilon i)=+1$, the born norm classes are
Galois-rational of balanced type by the transport law, and the R47/R48
identities are the $f=1$ instances. \textbf{Theorem 4 (the quaternion
deck)}: at Albert types II/III the four-power deck
$\gamma^2=-1$ lives \emph{inside} the endomorphism algebra, and the
alternating law $\trd(\gamma\,x\bar x)=\Nrd(x)\trd(\gamma)=0$ is the
grade-$2$ vanishing avatar. \textbf{Theorem 5 (the derived landing
cycles)}: the conduit's algebraic core (the Newton conductor, the twist
cancellation, the Thom--Porteous assembly) is proved and certified, and
the two seeds land \emph{derived} cycles --- the $\mathbf P^2$ pair at
$p=1$ and the committed $\Gr(2,5)$ pair $\xi_{360}$ at the
open-sector grade $p=2$, whose cycle
$W=\tfrac25\sigma_2-\tfrac35\sigma_{1,1}$ is computed in exact
arithmetic at three independent twists $L=\Oo(1),\Oo(2),\Oo(3)$. \textbf{Theorem 6 (the weld and
the FS-lock)}: the remaining falsifiable line for Q2 is now the
\emph{single} law FS --- the pair-realization
$m^*(X,p,\alpha)>0$ --- whose equivalence to the Clay statement is the
corpus's own certified receipt; the false pair $(C1,C2)$ is dissolved:
$C2$ dead, $C1$'s conduit a theorem at the seeds, the input axis
carrying the generalized fold-birth at every non-residue stratum.
\end{abstract}

\section{The state and the corpus}

Round 48's Definition 3 named the two commitments:
\begin{itemize}
\item \textbf{(C1)} the landing chain: the committed field
\texttt{FoldProcessOn.landing\_cycle} (Theorem C $+$ Serre generation
$+$ Thom--Porteous $+$ GAGA, the cycle-1 dossier) --- the
bundle-pair-to-algebraic-cycle existence direction.
\item \textbf{(C2)} the instantiation hypothesis: every smooth
projective complex variety's extracted input (the R32 middle
extraction $M_{\mathrm{mid}}\subseteq H^{2p-1}(X,\Q)$) is a
two-letter world --- carries integral $(s,q)$ letter parameters, so its
phase is quadratic over $\Q$.
\end{itemize}
The complete branch corpus (the $741$-file map of this round's opening
sweep: $729$ root $+$ $9$ wave $+$ $3$ challenge Lean files, zero
\texttt{axiom} declarations, zero \texttt{native\_decide}, the CI-green
capstone \texttt{GSTGraphV2OntologicalFoldForAllN.lean} with
\texttt{FoldPairCalculus}, the proven scale-group principality
\texttt{scaleSet\_principal}, and the draft LawStack rounds with their
external certificates) was mapped in full by a three-agent dispatch
this round: a corpus map, a derivation extraction, and an independent
war council. The council's verdict, which this round now executes:
$C1$ and $C2$ are \emph{one wall photographed from two sides} (the
missing pair-realization $\Sigma$), plus \emph{one false claim}. The
false claim dies here (Theorem 1); the wall is named once (Theorem 6);
and the fold equations --- the operator's main idea --- are generalized
to the entire arithmetic input space (Theorems 2--4).

\section{The death certificate}

\begin{theorem}[The death certificate for (C2)]
\label{th:death}
The instantiation hypothesis (C2) --- ``every smooth projective complex
variety's extracted input is a two-letter world'' --- is refuted as a
blanket modeling claim:
\begin{enumerate}[label=(\roman*)]
\item \textbf{The empty referent at the flagship.} Let $X$ be a smooth
cubic fourfold and $p=2$. Then $H^3(X,\Q)=0$, hence
$M_{\mathrm{mid}}(X,2)=0$: the extracted input is \emph{empty} and no
letter parameters exist, while $H^4(X,\Q)\cap H^{2,2}$ carries the
famous open Hodge-conjecture content of the cubic fourfolds. The weld's
``every extracted phase lands in a receipted stratum'' has no referent
at this input.
\item \textbf{A non-empty letterless input.} Let $A_1,A_2$ be very
general, non-isogenous polarized abelian surfaces and $X=A_1\times
A_2$, $p=2$. Then
\[
M_{\mathrm{mid}} \;=\; H^1(A_1)\otimes \Q\lambda_2 \;\oplus\;
\Q\lambda_1\otimes H^1(A_2) \;\neq\; 0,
\qquad \lambda_i := c_1(L_i),
\]
with types $\subseteq\{(1,2),(2,1)\}$, and
\[
\End_{\Hdg}(M_{\mathrm{mid}}) \;=\; \Q\oplus\Q,
\]
which admits \emph{no} quadratic letter: there is no injective unital
$\Q$-algebra map $\Q(\sqrt D)\to\Q\oplus\Q$ for any non-square $D$.
\item \textbf{Meagreness.} Within the polarized abelian input space,
the lettered inputs form a meagre subset: for each fixed integral
letter candidate the Hodge-endomorphism condition is a proper analytic
condition on the period matrix, and there are countably many
candidates. By Baire the complement is dense.
\end{enumerate}
Consequently (C2) is not a hypothesis to be proved but a modeling
error to be buried; its replacement is the Albert stratification of
Theorem~\ref{th:albert}.
\end{theorem}

\begin{proof}
(i) The cubic fourfold is a smooth hypersurface in $\mathbf P^5$. By
the Lefschetz hyperplane theorem the inclusion
$i\colon X\hookrightarrow\mathbf P^5$ induces isomorphisms
$H^k(\mathbf P^5,\Q)\to H^k(X,\Q)$ for $k<\dim X=4$. Since
$H^3(\mathbf P^5,\Q)=0$ (the cohomology of projective space is
$\Q[h]/(h^6)$, concentrated in even degrees), $H^3(X,\Q)=0$ and the
maximal level-$\le 1$ rational sub-Hodge-structure of $H^3$ is $0$.
That the rational $(2,2)$ classes on cubic fourfolds are the standard
open instance of the Hodge conjecture (the special cubics of Hassett's
program) is classical. (ii) By K\"unneth,
$H^3(A_1\times A_2)=H^1(A_1)\otimes H^2(A_2)\oplus H^2(A_1)\otimes
H^1(A_2)$. The level of a tensor product of Hodge structures is the
sum of the levels; the level-$\le 1$ part therefore comes from
$H^1\otimes(\text{level }0)$, and the level-$0$ rational part of
$H^2(A_i)$ is $H^2(A_i,\Q)\cap H^{1,1}=\NS(A_i)\otimes\Q=\Q\lambda_i$
by Lefschetz $(1,1)$ and the very generality of the polarization
($\rho=1$). The types of $H^1\otimes\Q\lambda$ are
$\{(1,0)+(1,1),(0,1)+(1,1)\}=\{(2,1),(1,2)\}$ as claimed. For the
endomorphisms: $\End_{\Hdg}$ respects the two tensor factors, and
\[
\End_{\Hdg}(M_{\mathrm{mid}})=
\End_{\Hdg}(H^1A_1)\oplus\End_{\Hdg}(H^1A_2)\oplus
\Hom_{\Hdg}(H^1A_2,H^1A_1)\oplus\Hom_{\Hdg}(H^1A_1,H^1A_2).
\]
The cross terms are Hodge classes in $H^1(A_i)^\vee\otimes H^1(A_j)$,
which vanish: the classical Hodge-correspondence theorem for abelian
varieties gives $\Hom_{\Hdg}(H^1(A_i),H^1(A_j))=\Hom(A_i,A_j)\otimes\Q$,
and $\Hom(A_1,A_2)=0$ for non-isogenous very general abelian
surfaces. The diagonal terms are
$\End_{\Hdg}(H^1(A_i))=\End(A_i)\otimes\Q=\Q$ by very generality. Hence
$\End_{\Hdg}(M_{\mathrm{mid}})=\Q\oplus\Q$. Finally, if
$\varphi\colon\Q(\sqrt D)\to\Q\oplus\Q$ were an injective unital
algebra map, then $\varphi(\sqrt D)=(a,b)$ with $a^2=b^2=D$,
$a,b\in\Q$ --- impossible for non-square $D$; equivalently, a field has
no nontrivial idempotents while $\Q\oplus\Q$ has $(1,0)$.
(iii) A letter on $M_{\mathrm{mid}}(A,p)$ is in particular a rational
endomorphism of the lattice $H_{2p-1}(A,\Z)$ of trace $s$ and
determinant $q$ preserving the Hodge decomposition. For a fixed
rational matrix $R$ the condition ``$R$ preserves the Hodge
decomposition'' is a system of linear equations in the period matrix
$Z$ of the polarized abelian variety; it defines a closed analytic
subset of the Siegel domain, proper unless $R$ preserves every Hodge
decomposition, which for non-scalar $R$ fails at a very general period
(each $\End_{\Hdg}=\Q$ by the Mattuck--Moonen--Zarhin triviality of
the Hodge ring at the very general point). Countably many $R$; Baire
applies on the complete metric space of periods.
\end{proof}

\section{The Albert stratification}

\begin{theorem}[The Albert stratification: the true world model]
\label{th:albert}
Fix a smooth projective complex $X$ and $p$ with
$M_{\mathrm{mid}}(X,p)\neq 0$. Then $J_{\mathrm{mid}}(X,p)$ is an
abelian variety (the R32 receipt) and
$D:=\End^0(J_{\mathrm{mid}})$ is a finite-dimensional semisimple
$\Q$-algebra carrying the positive Rosati involution. The input space
$\{(X,p):M_{\mathrm{mid}}\neq 0\}$ partitions into the Albert cells:
\[
\mathcal S_{\Q}\ (\text{trivial}),\quad
\mathcal S_{\mathrm I}\ (\text{totally real fields}),\quad
\mathcal S_{\mathrm{II}}\ (\text{indefinite quaternions}),\quad
\mathcal S_{\mathrm{III}}\ (\text{definite quaternions}),\quad
\mathcal S_{\mathrm{IV}}\ (\text{CM type}),
\]
by the simple components of $D$ (with their centers and involution
types, per Albert's classification). The transfer holds:
\begin{enumerate}[label=(\alph*)]
\item the R47 real-multiplication stratum is exactly the degree-$2$
cell of $\mathcal S_{\mathrm I}$; the R48 unit-circle stratum and the
committed CM censuses are degree-$2$ cells of $\mathcal S_{\mathrm{IV}}$
(a CM field has even degree);
\item every non-residue cell carries a number field $K\subseteq D$ of
even degree: type I/IV centers, or a quadratic subfield of a
quaternion component (at center $\Q$: the pure-quaternion square
roots); odd-degree totally real centers carry the odd-degree
obstruction of Remark~\ref{rmk:odd};
\item at every such $K$ the orbit law (Theorem~\ref{th:orbit}) fires;
at types II/III the quaternion deck (Theorem~\ref{th:quaternion}) fires.
\end{enumerate}
\end{theorem}

\begin{proof}
$J_{\mathrm{mid}}$ abelian: the R32 Theorem 1 receipt (the middle
extraction is two-level polarizable; its intermediate Jacobian is an
abelian variety). $\End^0$ of an abelian variety is semisimple with a
positive involution (the Rosati of any polarization), and its simple
components, with center a number field and involution of positive
type, are classified by Albert's four types --- classical (Albert;
cf. Birkenhake--Lange, Ch.~5; Shimura--Taniyama). The partition by the
multiset of component types is exhaustive and mutually exclusive by
definition. (a) A real quadratic field acting on $J_{\mathrm{mid}}$ is
a degree-$2$ totally real subalgebra, i.e.\ a type I cell of degree
$2$; an imaginary quadratic (UC/census) action is a degree-$2$ CM
subfield, i.e.\ a type IV cell of degree $2$. (b) A CM field has even
degree (complex conjugation acts fixed-point-freely on its
embeddings). A quaternion algebra over $\Q$ contains the quadratic
field $\Q(\sqrt a)$ (the defining generator squares to $a$); over a
totally real center $E$ of degree $e$, the component contains maximal
commutative subfields of degree $2e$, and the even-degree letter is
supplied by $E$ itself when $e$ is even. (c) is Theorems
\ref{th:orbit} and \ref{th:quaternion}.
\end{proof}

\begin{remark}[The odd-degree obstruction]
\label{rmk:odd}
For $K$ totally real of odd degree $e$, the norm product of the cells
has type $(ep-k, ep-e+k)$ with $k\neq e/2$ impossible: the balanced
$(q,q)$ type cannot be reached. The letter method's honest boundary:
odd-degree totally real centers feed the residue line (Theorem
\ref{th:weld}), not the gate.
\end{remark}

\section{The orbit law}

\begin{theorem}[The orbit law: the generalized phase-pair cancellation]
\label{th:orbit}
Let $M$ be a polarizable two-level $\Q$-Hodge structure of types
$\{(p-1,p),(p,p-1)\}$ (the middle extraction), carrying an action of a
number field $K\subseteq\End_{\Hdg}(M)$ of even degree $n=2f$. Write
$M\otimes\C=\bigoplus_{\varphi\in\Emb(K)}M_\varphi$ for the eigencell
decomposition. Then:
\begin{enumerate}[label=(\roman*)]
\item \textbf{The transport law.} Complex conjugation $c$ on
$M\otimes\C$ maps $M_\varphi$ to $M_{\bar\varphi}$ and exchanges the
Hodge types: $c(M^{a,b})=M^{b,a}$; hence $c$ pairs each cell of type
$(p-1,p)$ with a cell of type $(p,p-1)$, and splits each
real-embedding cell internally into the two types in equal dimension.
\item \textbf{The balance (fold equation 1).}
$\sum_\varphi\dim(M_\varphi\cap M^{p-1,p})=\sum_\varphi
\dim(M_\varphi\cap M^{p,p-1})=\dim M/2$.
\item \textbf{The unified cancellation (the $(p,p)$-gate).} On a cell
of type $(p,p-1)$ the Weil operator acts by $+i$; on a cell of type
$(p-1,p)$ by $-i$. Every conjugate pair of cells therefore contributes
the eigenvalue identity
\[
(\varepsilon\, i)\bigl(-\varepsilon\, i\bigr) \;=\; +1,
\qquad \varepsilon=\pm1,
\]
and the Galois-norm orbit product of a cell family carries Weil
eigenvalue $+1$ at \emph{every} even degree $n=2f$, real or imaginary
embeddings, $f$ arbitrary. The R47/R48 phase-pair cancellations are the
$f=1$ instances.
\item \textbf{The born classes.} The symmetrized orbit products
$\mathrm{Sym}\bigl(\bigotimes_\varphi x_\varphi\bigr)$ of cell
families are $\Gal$-invariant, hence defined over $\Q$ (descent), and
of balanced type $(f(2p-1),f(2p-1))$; with the Gaussian shift
$H^1(E)\otimes M_\subseteq H^{2p}(E\times X)$ the cells are
$(p,p)$-typed and the born classes are $(np,np)$-Hodge classes on the
product. \emph{The fold-birth: Hodge classes born from the
transcendental cell data by the conjugation pairing.}
\item \textbf{The projector spread.} For every
$\Gal$-stable $S\subseteq\Emb(K)$ the spectral sum
$P_S:=\sum_{\varphi\in S}e_\varphi\in(K\otimes\C)^{c}=K$ is a
$\Q$-rational Hodge endomorphism with $P_S^2=P_S$ and
$P_S+P_{S^c}=\mathrm{id}$. The R48 projector law
$P_A+\sigma(P_A)=I$ is the $|S|=1$, $n=2$ instance.
\item \textbf{The five fold equations at the orbit level.}
(1) the balance is (ii); (2) the grade-$2$ vanishing: for every
anti-invariant $v$ ($\sigma(v)=-v$, with $\sigma$ an involution of
$K$ --- present at every even-degree Galois letter and every CM
letter; odd-degree letters admit none, consistent with
Remark~\ref{rmk:odd}) the trace pairing
$\Psi_v(x,y):=\tr_{K/\Q}(v\,x\,\sigma(y))$ is \emph{alternating} ---
$\Psi_v(y,x)=-\Psi_v(x,y)$ at every Galois degree; (3) the trace law
$\ch(P)-\ch(M)=m\alpha$ pure is the seed-side content of Theorem
\ref{th:derived} (the Newton conductor); (4) the scale $m\neq0$ is the
capstone's principality receipt (\texttt{scaleSet\_principal},
CI-green); (5) the rationality $N_{K/\Q}(x)\in\Q$ is the field norm.
\end{enumerate}
\end{theorem}

\begin{proof}
(i) $K$ acts by $\Q$-rational endomorphisms, so $c\,\theta\,c^{-1}
=\theta$ for $\theta\in K$; if $v\in M_\varphi$ (i.e.\
$\theta v=\varphi(\theta)v$ after choosing a primitive
$\theta$) then $\theta(cv)=c(\theta v)=
\overline{\varphi(\theta)}\,cv=\bar\varphi(\theta)\,cv$, so
$cM_\varphi=M_{\bar\varphi}$. And $c(M^{a,b})=M^{b,a}$ is standard for
a $\Q$-Hodge structure. For a real embedding $\bar\varphi=\varphi$: the
cell is $c$-stable, and $c$ exchanges its two type-parts, which
therefore have equal dimension. (ii) $c$ is an involution exchanging
$M^{p-1,p}$ and $M^{p,p-1}$ globally (bijective), and within split
cells equally; counting through the cell decomposition gives the
balance. (iii) The Weil operator acts on $H^{a,b}$ by $i^{a-b}$: on
$(p,p-1)$ by $i$, on $(p-1,p)$ by $-i$. A conjugate pair of cells
carries opposite types by (i), so the pair product --- the cup/tensor
product of one class from each --- has $C$-eigenvalue
$(\varepsilon i)\cdot(-\varepsilon i)=\varepsilon^2 i^2
\cdot(-1)\cdot(-1)\cdot$~$\ldots$~$=+1$ (the sign identity
$i\cdot(-i)=1$ with $\varepsilon^2=1$). Real cells contribute their
internal pairs with the same identity. The product over the orbit is
$(+1)^f=+1$. (iv) $\Gal$-invariance of the symmetrized products is
immediate (the Galois action permutes cells and tensor slots);
descent: the invariants of $V\otimes\bar\Q$ under $\Gal(\bar\Q/\Q)$
recover the $\Q$-structure. The type of a product of cell classes is
$(\sum a_i,\sum b_i)$ with the multiset $\{a_i\}$ containing $k=n/2$
copies of $p-1$ and $n/2$ copies of $p$ (the balance), giving
$(np-n/2,\,np-n/2)=(f(2p-1),f(2p-1))$. With the shift by
$H^1(E)=(1,0)\oplus(0,1)$: $(1,0)\otimes(p-1,p)$ and
$(0,1)\otimes(p,p-1)$ both have type $(p,p)$, so each shifted cell is
$(p,p)$-typed and the orbit products are $(np,np)$-typed. (v) The
idempotents $e_\varphi$ of $K\otimes\C\simeq\bigoplus_\varphi\C$ sum
over $\Gal$-stable $S$ to a $c$-invariant element, and
$(K\otimes\C)^c=K$ (an element of $K\otimes\C$ fixed by the
conjugation lies in $K\otimes(\C^c\cap$ rational diagonal$)$; for the
$\Q$-algebra $K\otimes_\Q\C$ the $c$-invariants are exactly $K$).
Orthogonality gives idempotence and the complement law.
(vi)(2) Galois case: write the trace as the embedding sum and
reindex by $g\mapsto g\circ\sigma$:
\[
\tr(v\,y\,\sigma x)=\sum_g g(v)g(y)(g\sigma)(x)
=\sum_{g'}(g'\sigma^{-1})(v)\,(g'\sigma^{-1})(y)\,g'(x)
=\tr(\sigma v\,\sigma y\,x)=-\tr(v\,x\,\sigma y),
\]
using commutativity, cyclicity, and $\sigma v=-v$ (with
$\sigma=\sigma^{-1}$). For non-Galois $K$ the same identity descends
from the Galois closure. The quadratic instances are the R48
$\psi(x,y)=\tr(\sqrt{D_2}\,x\,\sigma y)$ alternating law and the R44
isotropic Gram receipts. (vi)(5) $N_{K/\Q}(x)=x\,\sigma(x)\cdot(\text{orbit
factors})\in\Q$ is the field norm; the R48 identity $z\sigma(z)=1$ is
the unit-circle instance.
\end{proof}

\section{The quaternion deck}

\begin{theorem}[The quaternion deck at Albert types II/III]
\label{th:quaternion}
Let $D\subseteq\End^0(J_{\mathrm{mid}})$ be a quaternion component
over $\Q$ (type II: indefinite, splitting at $\infty$; type III:
definite, ramifying at $\infty$), with the Rosati involution
$'$ and reduced trace/norm $\trd,\Nrd$.
\begin{enumerate}[label=(\roman*)]
\item \textbf{The deck.} Type II: $D$ contains an element $\gamma$
with $\gamma^2=-1$ (e.g.\ $D=(a,b)$ with $a<0<b$: the generator
squares to $a$; rescale by $1/\sqrt{|a|}$). Type III (e.g.\ the
Hamilton algebra $(-1,-1)$): every pure quaternion
$\gamma=xi+yj+zij$ has
\[
\gamma^2=ax^2+by^2-ab\,z^2\big|_{a=b=-1}=-(x^2+y^2+z^2)<0,
\]
so $\gamma^2=-d$, $d>0$, and $\gamma^4=d^2>0$. In both cases the
four-power deck $\gamma^4>0$ is an \emph{internal} endomorphism: the
R47/R48 Gaussian deck $[i]\times\mathrm{id}$ (external, on the
$E$-factor) is internalized at the quaternion strata.
\item \textbf{The alternating law (grade-$2$ vanishing avatar).} For
pure $\gamma$ ($\trd(\gamma)=0$, $\bar\gamma=-\gamma$):
\[
\Psi_\gamma(x,x)=\trd(\gamma\,x\bar x)=\Nrd(x)\,\trd(\gamma)=0,
\]
so the pairing $\Psi_\gamma(x,y)=\trd(\gamma x\bar y)$ is alternating
on the nose --- the quaternion instance of fold equation (2).
\item \textbf{The Rosati positivity.} $\trd(d\,d')>0$ for all $d\neq0$
(the Albert receipt); at the Hamilton algebra,
$\trd(d\bar d)=2\,\Nrd(d)=2(w^2+x^2+y^2+z^2)>0$ exactly.
\item \textbf{The bridge.} $D$ contains the CM quadratic field
$\Q(\sqrt a)$ (the generator's square root), so the $f=1$ orbit law
fires inside the quaternion strata; the deck $\gamma$ itself obeys the
transport law (its eigenvalues $\pm i$ pair under conjugation with the
type flip), giving the $C_4$-cell structure of the orbit law with
$\mathrm{Gal}=\langle\gamma\rangle$.
\end{enumerate}
\end{theorem}

\begin{proof}
(i) is the explicit quaternion arithmetic
$(xi+yj+zij)^2=ax^2+by^2-abz^2$ (the mixed terms cancel by the
anticommutation $ij=-ji$), verified in the certificate over the
definite and indefinite families. (ii) For $q\in\Q$ the reduced trace is $\trd(q\cdot1)=2q$
(the scalar acts on the $2$-dimensional algebra), while every pure
quaternion has $\trd(\gamma)=0$ (the reduced trace reads the scalar
part). With $x\bar x=\Nrd(x)\cdot1\in\Q$:
\[
\Psi_\gamma(x,x)=\trd\bigl(\gamma\,\Nrd(x)\bigr)
=\Nrd(x)\,\trd(\gamma)=0,
\]
so $\Psi_\gamma$ vanishes on the diagonal --- alternating on the nose. (iii) is Albert's
positivity theorem; the Hamilton computation is exact. (iv) The
generator $i$ of $(a,b)$ satisfies $i^2=a$: for $a<0$ non-square,
$\Q(i)=\Q(\sqrt a)$ is CM. The transport law for $\gamma$ is the same
conjugation-commutation argument as in Theorem~\ref{th:orbit}(i).
\end{proof}

\begin{remark}
The certificate verifies the trace statements on the exact
$4$-dimensional algebras: $\trd(q\cdot1)=2q$ on scalars,
$\trd(\gamma)=0$ on pure quaternions, at the definite and indefinite
families.
\end{remark}

\section{The derived landing cycles}

\begin{lemma}[The conduit's algebraic core]
\label{lem:conduit}
Let $X$ be smooth projective, $E_\pm$ holomorphic of equal rank $r$,
with $\ch(E_+)-\ch(E_-)=m\alpha$ \emph{pure} at grade $p$ ($m\neq0$),
and let $L$ be any line bundle, $\ell=c_1(L)$. Then:
\begin{enumerate}[label=(\alph*)]
\item \textbf{The Newton conductor.} For the virtual $\xi:=E_+-E_-$
(with $c(\xi):=c(E_+)c(E_-)^{-1}$ and virtual roots), the power sums
$p_j(\xi)=j!\,\ch_j(\xi)$ vanish for $j<p$ and $p_p(\xi)=p!\,m\alpha$;
Newton's identities give $e_j(\xi)=0$ for $j<p$ and
\[
c_p(\xi)\;=\;e_p(\xi)\;=\;(-1)^{p-1}(p-1)!\,m\,\alpha.
\]
\item \textbf{The twist cancellation.} The same holds for
$\xi\otimes L$: its lower power sums vanish \emph{regardless of
$\ell$} (only grades $\ge p$ of $\ch(\xi)e^\ell$ are nonzero below
twisting by $\ell$ appears first at grade $p{+}1$), so
\[
c_p(E_+\otimes L)-c_p(E_-\otimes L)\;=\;c_p(\xi\otimes L)
\;=\;(-1)^{p-1}(p-1)!\,m\,\alpha
\]
\emph{at every twist}, the difference being twist-invariant.
\item \textbf{The Thom--Porteous assembly.} For $s=r-p+1$ general
sections $\sigma_\pm\colon\Oo^s\to E_\pm\otimes L$ with determinantal
centers $D_\pm=D_{s-1}(\sigma_\pm)$ (expected codimension $p$), the
Porteous class is $[D_\pm]=c_p(E_\pm\otimes L)$ (the $1\times1$
determinant of the Thom--Porteous formula; the source is trivial), so
\[
[D_+]-[D_-]\;=\;(-1)^{p-1}(p-1)!\,m\,\alpha,
\qquad
W:=\tfrac{(-1)^{p-1}}{(p-1)!\,m}\bigl([D_+]-[D_-]\bigr),\quad
\Cl(W)=\alpha.
\]
\end{enumerate}
\end{lemma}

\begin{proof}
(a) The virtual power sums are $p_j(\xi)=j!\ch_j(\xi)$ because
$\ch=\sum_i e^{x_i}$ over the virtual multiset of roots. Newton's
identity $ke_k=\sum_{i=1}^k(-1)^{i-1}e_{k-i}p_i$ is triangular: if
$p_i=0$ for $i<p$ then $e_j=0$ for $j<p$ inductively, and at $k=p$
only the $i=p$ term survives:
$pe_p=(-1)^{p-1}e_0p_p$, i.e.\
$e_p=(-1)^{p-1}(p-1)!\,m\alpha$. (b)
$\ch(\xi\otimes L)=\ch(\xi)e^\ell$: its grade-$j$ part is
$\sum_{a\le j}\ch_a(\xi)\ell^{j-a}/(j-a)!=0$ for $j<p$ (all
$\ch_a(\xi)=0$ for $a<p$, including $a=0$ since the ranks are equal),
and $=m\alpha$ at $j=p$ (the single surviving term $a=p$). The
conductor of (a) applies verbatim to $\xi\otimes L$. The identity
$c_p(E_+\otimes L)-c_p(E_-\otimes L)=c_p(\xi\otimes L)$: expand
$c(E_+\otimes L)=c(\xi\otimes L)\,c(E_-\otimes L)$ (the definition of
the virtual total Chern class is exactly this multiplicativity);
comparing grade $p$, every cross term contains a factor
$e_j(\xi\otimes L)$ with $0<j<p$, which vanishes. (c) For a morphism
$\sigma\colon\Oo^s\to F$ of rank $r$ with $s=r-p+1$ general sections,
$D_{s-1}(\sigma)$ has expected codimension $(r-s+1)\cdot1=p$ and
Thom--Porteous gives
$[D_{s-1}]=\Delta_{(p)}\bigl(c(F)\bigr)=c_p(F)$ --- the $1\times1$
determinant (Fulton, \emph{Intersection Theory}, Thm.~14.4(a); the
source is trivial of rank $s$, so $c(F-\Oo^s)=c(F)$). GAGA: the
determinantal centers of algebraic sections are algebraic cycles. The
assembly is (b) plus (c).
\end{proof}

\begin{theorem}[The derived landing cycles at the two seeds]
\label{th:derived}
\begin{enumerate}[label=(\roman*)]
\item \textbf{The $\mathbf P^2$ seed at $p=1$.} The pair
\[
E_+=\Oo(1)\oplus\Oo(1),\qquad E_-=\Oo(1)\oplus\Oo(-1)
\]
has equal ranks $2$, and
$\ch(E_+)-\ch(E_-)=\ch(\Oo(1))-\ch(\Oo(-1))=2h$ \emph{pure at grade
$1$} ($m=2$, $\alpha_1=h$). The conduit lands
\[
[D_+]-[D_-]=c_1\text{-difference}=2h,\qquad
W=\tfrac12([D_+]-[D_-])\sim \text{hyperplane},\qquad \Cl(W)=h:
\]
the derived landing cycle at the divisor grade, cross-receipted by
Lefschetz $(1,1)$.
\item \textbf{The $\Gr(2,5)$ seed at the open-sector grade $p=2$.}
The gr25 committed class
\[
\xi_{360}=351\,B_{(1,0)}+5\,B_{(3,0)}+736\,B_{(1,1)}+36\,B_{(0,2)}
+B_{(0,3)}
-17\,B_{(0,0)}-51\,B_{(2,0)}-1620\,B_{(0,1)}-141\,B_{(2,1)}-9\,B_{(1,2)}
\]
($B_{(j,i)}=\Sym^jS^\vee\otimes\Oo(i)$; minimal scale $m_0=360$)
splits into the equal-rank pair
\[
E_+ = 351\,S^\vee+5\,\Sym^3S^\vee+736\,S^\vee(1)+36\,\Oo(2)+\Oo(3),
\qquad \mathrm{rk}=2231,
\]
\[
E_- = 17\,\Oo+51\,\Sym^2S^\vee+1620\,\Oo(1)+141\,\Sym^2S^\vee(1)
+9\,S^\vee(2),\qquad \mathrm{rk}=2231,
\]
with $\ch(E_+)-\ch(E_-)=360\alpha$, $\alpha=\tfrac25h^2-s_2$, pure at
grade $2$. The conduit lands, at \emph{every} twist $L$:
\[
[D_+]-[D_-]=c_2\text{-difference}=-360\,\alpha,\qquad
W=\tfrac{[D_-]-[D_+]}{360}=\alpha=\tfrac25\,\sigma_2-\tfrac35\,\sigma_{1,1},
\]
an explicit $\Q$-combination of the two codimension-$2$ Schubert
varieties. The certificate verifies the ch-purity, the equal ranks,
the twist-invariance at $L=\Oo(1)$, $\Oo(2)$ and $\Oo(3)$, and the Schubert
landing, all in exact arithmetic.
\end{enumerate}
In consequence the \texttt{landing\_cycle} field's monopoly is broken:
at both seeds the cycle is \emph{derived} from the world's own
$K^0$-data by the certified conduit --- the first derived landing
cycles of the campaign, at the divisor grade \emph{and} at the
open-sector grade $p=2$.
\end{theorem}

\begin{proof}
(i) The ch-computation on $\mathbf P^2$:
$\ch(\Oo(k))=1+kh+\tfrac{k^2h^2}{2}$ (the series truncates at $h^2$),
so $\ch(\Oo(1))-\ch(\Oo(-1))=2h$ exactly. The conduit is Lemma
\ref{lem:conduit} at $p=1$. (ii) The purity and the coordinates are
the gr25 theorem, re-verified independently by this round's
certificate against the corpus component table (the $10\times10$
Schur ch-table reproduced exactly). The equal ranks follow from
$\ch_0(\xi_{360})=0$. The pair split is the positive/negative part.
Lemma \ref{lem:conduit}(b)--(c) at $p=2$ gives the Porteous
difference $(-1)^{1}1!\,360\alpha=-360\alpha$ at every twist; the
certificate computes $c_2(E_\pm\otimes\Oo(t))$ by exact Newton
inversion of $\ch(E_\pm)e^{th}$ in the Schubert ring
$\Q[h,s_2]/(h^4-3h^2s_2+s_2^2,\;h^3s_2-2hs_2^2)$ at $t=1,2$ and
verifies the invariance. The Schubert translation
$h^2=\sigma_2+\sigma_{1,1}$, $s_2=\sigma_{1,1}$ (the corpus integral
receipt $\int s_2h^4=2=\int\sigma_{1,1}h^4$) gives
$\alpha=\tfrac25\sigma_2-\tfrac35\sigma_{1,1}$. The geometric
realization: the degeneracy loci of general sections of tautological
bundles on the Grassmannian are Schubert-type subvarieties; the class
identity is the certified Porteous computation.
\end{proof}

\section{The weld and the FS-lock}

\begin{theorem}[The weld and the FS-lock]
\label{th:weld}
The state of the campaign after this round:
\begin{enumerate}[label=(\roman*)]
\item \textbf{The certified scope.} Every rational $(p,p)$ class is
algebraic on: the $114$ committed worlds (R46); the seeds' worlds
($\mathbf P^2$ and $\Gr(2,5)$, by Theorem \ref{th:derived}'s derived
landings); the frame's mixed sector at its committed scope (R32/R41,
the AJ-correspondence equivalence). The input axis of the true input
space carries the generalized fold-birth at every non-residue Albert
cell (Theorems \ref{th:albert}--\ref{th:quaternion}): the five fold
equations, the $(p,p)$-gate, the projector spread, at every even-degree
letter.
\item \textbf{The FS-lock.} The remaining falsifiable line for Q2 ---
the Clay statement at the full generality of every smooth projective
variety --- is the \emph{single} law
\[
\mathrm{FS}:\qquad
m^*(X,p,\alpha)>0\ \ \text{for every rational }(p,p)\text{ class:}
\ \ \exists\ \text{an integrable pure pair }(E_+,E_-)\text{ on }X
\text{ with }\ch(E_+)-\ch(E_-)=m\alpha,\ m\neq0,
\]
together with its letter-world avatar (the realization of the
cohomological norm pairs of Theorem \ref{th:orbit} by actual
$K^0$-classes). By the corpus's certified packet-equivalence (the gr25
non-circularity remark: the input is equivalent in strength to the
Hodge conjecture, both directions; the R32 remark), FS $\iff$ Q2.
\item \textbf{The dissolution of the pair.} $(C1,C2)$ is dissolved:
$C2$ is dead (Theorem \ref{th:death}) and replaced by the true
stratification; $C1$'s conduit is a theorem at the seeds (Theorem
\ref{th:derived}); the one wall --- the pair-realization $\Sigma$ ---
is named once, as FS.
\end{enumerate}
\end{theorem}

\begin{proof}
(i) is the corpus's receipts plus this round's Theorems
\ref{th:derived} and \ref{th:albert}--\ref{th:quaternion}. (ii) The
equivalence receipts are the gr25 dossier's non-circularity remark and
the R32 remark, both certified in prior rounds: a pure pair at a
rational $(p,p)$ class exists iff the class is algebraic (the forward
direction by the conduit, Lemma \ref{lem:conduit}; the backward
direction by resolving algebraic cycles into pairs). (iii) collects
Theorems \ref{th:death} and \ref{th:derived}.
\end{proof}

\section{The certificate}

\texttt{albert\_fold\_birth\_round49\_checks.py}: exact arithmetic
(\texttt{Fraction} only, zero floats), deterministic double-run, md5
seal, timer-stripped, with layers:
\begin{itemize}
\item \textbf{X0} environment: the float-scan, the import discipline.
\item \textbf{X1} the death certificate: the cubic-fourfold Betti
bookkeeping ($H^3(\mathbf P^5)=0\Rightarrow H^3(X)=0$ by Lefschetz
hyperplane; $b_4=23$ with $h^{3,1}=h^{1,3}=1$, $h^{2,2}=21$); the
$A_1\times A_2$ K\"unneth/level decomposition; the non-embedding
$\Q(\sqrt D)\not\hookrightarrow\Q\oplus\Q$ for a family of $D$; the
countable-union meagreness model.
\item \textbf{X2} the Albert layer: the five cells; the CM even-degree
law; the type table (center, $\R$-realization, involution) per
Albert type; the degree table of the R47/R48 strata.
\item \textbf{X3} the orbit law: the pair identity
$(\varepsilon i)(-\varepsilon i)=+1$ over all sign patterns and all
even $n\le24$ (exact $\Z/4$ arithmetic on $i$-powers); the balanced
type arithmetic $(np-k,\,np-n+k)$ with $k=n/2$; the alternating trace
law $\tr(vy\sigma x)=-\tr(vx\sigma y)$ over quadratic, biquadratic
and quartic cyclotomic fields (exact $\Q$-bases); the norm
rationality $N_{K/\Q}(x)\in\Q$; the projector idempotence and orbit
sums over the same fields.
\item \textbf{X4} the quaternion layer: the algebras $(a,b)$ as exact
$4$-dimensional $\Q$-algebras; the pure-square law
$\gamma^2=ax^2+by^2-abz^2$; the Hamilton definiteness; the deck
$\gamma^2=-1$ at $(-1,2)$; the alternating law
$\trd(\gamma x\bar x)=0$; the Rosati positivity
$\trd(d\bar d)=2\Nrd(d)>0$ at the Hamilton algebra; the Hilbert
symbol ramification classification over the primes $\le 100$ (exact
quadratic-residue arithmetic).
\item \textbf{X5} the conduit: the Newton conductor verified by exact
symmetric-function arithmetic at $p\le8$ (power sums to elementary
symmetric functions, the triangular induction); the twist
cancellation at the instance level.
\item \textbf{X6} the seeds: the full $\Gr(2,5)$ computation of the
exploration (the ring, the $10$-row ch-table against the corpus, the
$\xi_{360}$ purity, the pair split, the three-twist Porteous
difference, the Schubert landing); the $\mathbf P^2$ pair.
\item \textbf{X7} the weld: the scope/receipt matrix and the FS-lock
composition.
\item \textbf{H1--H10} the hostile battery: sabotages (delete a
stratum, flip an orbit sign, rig the Betti table, fake the conductor,
corrupt the ring relations, swap the pair split, fake the twist,
instantiate a fake landing, break the count seal, insert a float) ---
each detected by a redundant independent layer.
\end{itemize}

\section{The audit protocol}

AUDIT-49: a fresh-agent dispatch (the fifth fully independent audit in
campaign history), armed with the adversarial-proof-audit checklist
(claim match, hidden strengthenings, circular reductions, uniformity,
definition drift, computational leaps, case leaks, citation accuracy,
hostile counterexamples, blocked-route smell). One loop: FAIL
$\to$ repair $\to$ PASS. The audit re-derives: the Lefschetz
hyperplane vanishing, the K\"unneth level decomposition, the
alternating trace identity, the quaternion arithmetic, the Newton
conductor, the two-twist Porteous difference, and runs its own
sabotage battery against the certificate.

\section{Scope discipline}

Four classes, stated plainly:
\begin{enumerate}[label=(\alph*)]
\item \emph{Standard mathematics, used at proven scope}: the Lefschetz
hyperplane theorem; K\"unneth and level bookkeeping; Lefschetz
$(1,1)$; the Weil operator's type action; Galois descent; Newton's
identities; the splitting principle; the Thom--Porteous formula
(Fulton Thm.~14.4(a)); GAGA; Serre generation; Albert's
classification; the Rosati positivity; Baire's theorem.
\item \emph{The campaign's certified receipts}: R32 (the middle
extraction, the input-axis decomposition); the censuses (R35--R41,
R43); R44/R45; R46 (the leaf discharge at all $114$ worlds); R47 (the
RM fold-birth); R48 (the UC fold-birth, the sign-blind law); the gr25
dossier (Theorem C, the $\xi_{360}$ lattice, the minimal scale
$360$); the CI-green capstone (the fold-pair calculus, the
principality \texttt{scaleSet\_principal}, \texttt{p2\_mstar\_eq\_two});
the packet-equivalence receipts.
\item \emph{This round's new mathematics}: the death certificate
(Theorem \ref{th:death}); the Albert stratification transfer
(Theorem \ref{th:albert}); the orbit law (Theorem \ref{th:orbit}); the
quaternion deck (Theorem \ref{th:quaternion}); the conduit's
assembled core (Lemma \ref{lem:conduit}); the derived landing cycles
(Theorem \ref{th:derived}); the FS-lock (Theorem \ref{th:weld}).
\item \emph{The commitment class}: FS --- the pair-realization, one
named law, equivalent in strength to the Clay statement by the
corpus's own receipts; the letter-world realization step (the norm
pairs of Theorem \ref{th:orbit} as $K^0$-classes) is the same wall
photographed from the letter side. No receipt is silently
strengthened; the closing string is emitted at exactly the scope of
(i)--(iii).
\end{enumerate}

The quantifier line: \textbf{the false hypothesis is dead; the true
input space is stratified; the fold equations hold at every even-degree
letter in the true stratification; the landing cycles are derived at
both seeds including the open-sector grade; the remaining falsifiable
line for Q2 is the single law FS.} The campaign continues per the
operator's standing order.

\end{document}
